\section{Korelace}\label{sec:ai-korelace}

Mějme dva stejně dlouhé statistické soubory
\[
    \mathbf{X}=(x_1,\ldots,x_n),
    \qquad
    \mathbf{Y}=(y_1,\ldots,y_n).
\]
Dvojice \((x_i,y_i)\) popisuje dva znaky téhož \(i\)-tého objektu. Bodový graf těchto dvojic často naznačí, zda se s~růstem jednoho znaku mění také druhý.

\begin{definition}\label{def:ai-kovariance}
    \emph{Kovariance} souborů \(\mathbf{X}\) a~\(\mathbf{Y}\) je číslo
    \[
        \Cov(\mathbf{X},\mathbf{Y})
        =
        \frac1n\sum_{i=1}^{n}
        (x_i-\average{\mathbf{X}})
        (y_i-\average{\mathbf{Y}}).
    \]
\end{definition}

Kladná kovariance obvykle znamená, že nadprůměrné hodnoty \(\mathbf{X}\) doprovázejí nadprůměrné hodnoty \(\mathbf{Y}\). Záporná kovariance odpovídá opačnému směru. Její velikost však závisí na jednotkách obou znaků, a~proto ji nelze přímo porovnávat mezi různými dvojicemi dat.

\begin{definition}\label{def:ai-korelace}
    Nechť \(\sigma_{\mathbf{X}}>0\) a~\(\sigma_{\mathbf{Y}}>0\).
    \emph{Korelační koeficient} souborů \(\mathbf{X},\mathbf{Y}\) definujeme vztahem
    \[
        \rho_{\mathbf{X}\mathbf{Y}}
        =
        \frac{\Cov(\mathbf{X},\mathbf{Y})}
        {\sigma_{\mathbf{X}}\sigma_{\mathbf{Y}}}.
    \]
\end{definition}

Ekvivalentně lze psát
\[
    \rho_{\mathbf{X}\mathbf{Y}}
    =
    \frac{
        \sum_{i=1}^{n}
        (x_i-\average{\mathbf{X}})
        (y_i-\average{\mathbf{Y}})
    }{
        \sqrt{\sum_{i=1}^{n}(x_i-\average{\mathbf{X}})^2}
        \sqrt{\sum_{i=1}^{n}(y_i-\average{\mathbf{Y}})^2}
    }.
\]
Korelace je tedy kovariance po odstranění měřítka obou znaků.

\begin{figure}[ht]
    \centering
    \begin{tikzpicture}[x=0.72cm,y=0.62cm,every node/.style={font=\normalsize}]
    \begin{scope}
        \draw[diredge] (0,0) -- (4.5,0);
        \draw[diredge] (0,0) -- (0,3.6);
        \foreach \x/\y in {0.5/0.5,1/0.9,1.5/1.1,2/1.7,2.5/2.0,3/2.6,3.5/2.8,4/3.3}
            \fill[blue!75!black] (\x,\y) circle (2pt);
        \node[below] at (2.2,-0.45) {\(\rho_{\mathbf{X}\mathbf{Y}}>0\)};
    \end{scope}
    \begin{scope}[xshift=6cm]
        \draw[diredge] (0,0) -- (4.5,0);
        \draw[diredge] (0,0) -- (0,3.6);
        \foreach \x/\y in {0.5/3.2,1/2.8,1.5/2.4,2/2.1,2.5/1.8,3/1.2,3.5/0.9,4/0.4}
            \fill[orange!85!black] (\x,\y) circle (2pt);
        \node[below] at (2.2,-0.45) {\(\rho_{\mathbf{X}\mathbf{Y}}<0\)};
    \end{scope}
    \begin{scope}[xshift=12cm]
        \draw[diredge] (0,0) -- (4.5,0);
        \draw[diredge] (0,0) -- (0,3.6);
        \foreach \x/\y in {0.5/2.3,1/0.7,1.5/3,2/1.6,2.5/2.7,3/0.5,3.5/2,4/1.1}
            \fill[green!55!black] (\x,\y) circle (2pt);
        \node[below] at (2.2,-0.45) {\(\rho_{\mathbf{X}\mathbf{Y}}\approx0\)};
    \end{scope}
\end{tikzpicture}

    \caption{Kladná, záporná a~nulová lineární korelace.}
    \label{fig:ai-korelace}
\end{figure}

\begin{proposition}\label{prop:ai-korelace-vlastnosti}
    Pro korelační koeficient souborů \(\mathbf{X},\mathbf{Y}\) platí:
    \begin{enumerate}
        \item \(-1\leqslant\rho_{\mathbf{X}\mathbf{Y}}\leqslant1\);
        \item \(\abs{\rho_{\mathbf{X}\mathbf{Y}}}=1\) právě tehdy, když
        \[
            \mathbf{Y}=\alpha\mathbf{X}+\beta
        \]
        pro vhodná \(\alpha,\beta\in\R\) s~\(\alpha\neq0\).
        Přitom \(\rho_{\mathbf{X}\mathbf{Y}}=1\) pro \(\alpha>0\) a
        \(\rho_{\mathbf{X}\mathbf{Y}}=-1\) pro \(\alpha<0\).
    \end{enumerate}
\end{proposition}

\begin{proof}
    Položme
    \[
        a_i=x_i-\average{\mathbf{X}},
        \qquad
        b_i=y_i-\average{\mathbf{Y}}.
    \]
    Podle Cauchyovy--Schwarzovy nerovnosti platí
    \[
        \abs{\sum_{i=1}^{n}a_ib_i}
        \leqslant
        \sqrt{\sum_{i=1}^{n}a_i^2}
        \sqrt{\sum_{i=1}^{n}b_i^2}.
    \]
    Protože oba výrazy pod odmocninou jsou kladné, můžeme jimi nerovnost vydělit a~dostaneme
    \[
        \abs{\rho_{\mathbf{X}\mathbf{Y}}}\leqslant1.
    \]

    Rovnost nastane právě tehdy, když jsou vektory odchylek lineárně závislé, tedy když existuje \(\alpha\neq0\) takové, že
    \[
        y_i-\average{\mathbf{Y}}
        =
        \alpha(x_i-\average{\mathbf{X}})
    \]
    pro všechna \(i\). Po úpravě
    \[
        y_i
        =
        \alpha x_i+
        \underbrace{
            \average{\mathbf{Y}}
            -\alpha\average{\mathbf{X}}
        }_{\beta},
    \]
    a~tedy \(\mathbf{Y}=\alpha\mathbf{X}+\beta\).
    Obráceně, z~této rovnosti ihned plyne lineární závislost obou vektorů odchylek, a~tedy rovnost v~Cauchyově--Schwarzově nerovnosti.

    Nakonec
    \[
        \Cov(\mathbf{X},\mathbf{Y})
        =
        \Cov(\mathbf{X},\alpha\mathbf{X}+\beta)
        =
        \alpha\Var(\mathbf{X}),
    \]
    takže znaménko korelace je znaménkem \(\alpha\).
\end{proof}

\warningbox{Korelace měří pouze lineární vztah. Hodnota \(\rho_{\mathbf{X}\mathbf{Y}}=0\) nevylučuje nelineární závislost. Korelace také sama o~sobě nedokazuje příčinnou souvislost; společný pohyb obou znaků může být způsoben třetí proměnnou nebo náhodou.}
